\chapter*{General Introduction}
\addcontentsline{toc}{chapter}{General Introduction}

The automotive industry is undergoing a significant technological transformation driven by the increasing integration of electronics, software, connectivity, and advanced communication systems into modern vehicles. As vehicle architectures become more complex, the volume of data exchanged between their different electronic components continues to increase. This evolution has created new communication requirements and has reinforced the importance of high-performance networking technologies, particularly in the context of Automotive Ethernet.

The growing complexity of automotive communication systems also creates important challenges for testing and validation activities. During these activities, network traffic can be captured and stored in the form of Packet Capture (PCAP) traces. These traces contain valuable information about the communications occurring within a network and can be used to investigate system behavior, protocol exchanges, and potential communication issues. However, a single test scenario may generate a large volume of captured data, making manual analysis increasingly time-consuming and difficult.

Moreover, captured network traces may contain repetitive, redundant, or low-value traffic in addition to information that is relevant to the analysis objective. Consequently, engineers may need to spend considerable time filtering and organizing captured communications before meaningful analysis can be performed. This highlights the importance of an efficient preprocessing stage capable of reducing unnecessary information while preserving the relevant data required for subsequent processing.

At the same time, the development of artificial intelligence techniques creates new opportunities for supporting network analysis and testing activities. AI-based systems can assist in the processing of large volumes of communication data and support tasks such as pattern recognition, event identification, traffic analysis, and anomaly detection. However, the effectiveness of such approaches depends strongly on the quality, relevance, and structure of the information provided to them. Preparing network traces in a consistent and lightweight form therefore represents an important prerequisite for efficient automated analysis.

It is within this context that the present internship project, entitled \textit{``PCAP Preprocessing AI Agent for Automotive Ethernet Testing''}, was carried out. The objective of the project is to develop an automated preprocessing approach for Automotive Ethernet communication traces. The proposed solution aims to process captured communication data, extract relevant information, filter unnecessary or redundant traffic, and generate structured outputs suitable for downstream processing and AI-based analysis.

This work is conducted within an industrial context related to automotive testing and validation. By automating part of the trace preparation process, the proposed approach seeks to reduce the amount of information that must be manually investigated and to provide a clearer and more structured communication context for subsequent analysis. The resulting preprocessing pipeline is intended to contribute to more efficient workflows for the investigation and preparation of network communication traces.

The remainder of this report is organized into five chapters. Chapter~1 presents the project context, including the host organization, the challenges associated with Automotive Ethernet testing, the identified limitations, and the proposed solution. Chapter~2 introduces the main technologies and tools relevant to the project, with particular attention to Automotive Ethernet, automotive communication protocols, and the technologies used for data processing and AI agent development. Chapter~3 presents the requirements and the overall architecture of the proposed system. Chapter~4 details the implementation of the solution, together with the testing, validation, and obtained results. Finally, Chapter~5 concludes the report by summarizing the main achievements, discussing the limitations of the work, and presenting possible directions for future improvements.